% Proof notes for the root article. Requires amsmath, amssymb, amsthm,
% booktabs, hyperref. The table input is supplied beside the certificate.
% This is a complete mathematical contribution, not a proposed proof.

\section{An exact proof of the rational tetralogarithm conjecture}
\label{rattetra:sec}

The manuscript's equation \texttt{golden:eq:tetra} reproduces the
six-term relation proposed by Vladimir Reshetnikov in 2015 and again
in 2017~\cite{Reshetnikov2015,Reshetnikov2017}. Its present proof discussion reduces it to a cross-base Kummer
specialization but does not prove that specialization. We fill this local
gap. The proof below is self-contained apart from elementary differential
identities for the classical polylogarithms and a small finite rational
certificate, supplied in full. It also proves a one-parameter functional
identity from which the rational evaluation follows. No claim of worldwide
priority for the underlying classical functional-equation mechanism is made.

\begin{theorem}[The six rational arguments]
\label{rattetra:thm:main}
Put $A=\log2$ and $B=\log3$. Then
\begin{align}
&720\operatorname{Li}_4(\tfrac12)
 -2160\operatorname{Li}_4(\tfrac13)
 +2160\operatorname{Li}_4(\tfrac23)
 +270\operatorname{Li}_4(\tfrac14)\notag\\
&\qquad+540\operatorname{Li}_4(\tfrac34)
 +135\operatorname{Li}_4(\tfrac19)\notag\\
&=19\pi^4+30(\pi^2-B^2)(10A^2-12AB+3B^2)
       -30A^2(19A^2-24AB+8B^2).
\label{rattetra:eq:main}
\end{align}
In particular, the intermediate relation \texttt{golden:eq:kummer}
in the manuscript is also proved, by the two elementary duplications
already recorded there.
\end{theorem}

\subsection{A real generalized Rogers function and its exact descent}

Write $\ell_n(x)=\operatorname{Re}\operatorname{Li}_n(x)$ on the
principal branch. For $x\in\mathbb R\setminus\{0,1\}$ set
$L=\log|x|$, $M=\log|1-x|$, so $\ell_1(x)=-M$. Define
\begin{align*}
R_2(x)&=\ell_2(x)-\tfrac12L\ell_1(x),\\
R_3(x)&=\ell_3(x)-L\ell_2(x)+\tfrac13L^2\ell_1(x),\\
R_4(x)&=\ell_4(x)-L\ell_3(x)
                +\tfrac12L^2\ell_2(x)-\tfrac18L^3\ell_1(x).
\end{align*}
These functions extend continuously by $R_n(0)=0$ and
$R_n(1)=\zeta(n)$. Direct telescoping differentiation gives
\begin{equation}
 dR_n(x)=\frac{(-1)^n(n-1)}{n!}
 L^{n-2}(L\,dM-M\,dL),\qquad n=2,3,4.
\label{rattetra:eq:derivative}
\end{equation}
At a crossing of $x=1$, all logarithmically singular terms have positive
powers of $L$, which tend to zero; this proves the stated continuity and
allows identities obtained on regular subintervals to be joined.

Let $V=\mathbb Q(t)^\times\otimes_{\mathbb Z}\mathbb Q$, written
additively, and put
\[
 \beta_n[f]=f^{\otimes(n-2)}\otimes(f\wedge(1-f)).
\]
The sign of a rational function is torsion and vanishes in $V$.
Rational prime factors do not vanish and must be retained.

\begin{lemma}[Descent with several independent logarithms]
\label{rattetra:lem:descent}
Let $f_i(t)$ be nonconstant rational functions with no zero or pole
on an interval $I$, and let $v_i\in\mathbb Q$ satisfy
$\sum_i v_i\beta_4[f_i]=0$. For every linear functional
$\eta:V\to\mathbb R$ and each $n=2,3,4$,
\[
 \sum_i v_i\eta(f_i)^{4-n}R_n(f_i(t))
\]
is constant on $I$, with the continuous interpretation at $f_i(t)=1$.
\end{lemma}

\begin{proof}
Contract $4-n$ of the first two tensor slots with $\eta$.
The result is $\sum_i v_i\eta(f_i)^{4-n}\beta_n[f_i]=0$.
Evaluate the remaining ordinary tensor slots by $g\mapsto\log|g(t)|$
and the wedge by
$g\wedge h\mapsto\log|g|\,d\log|h|-\log|h|\,d\log|g|$.
Equation~\eqref{rattetra:eq:derivative} gives a zero derivative on every
regular subinterval. Continuity joins the constants across the finitely
many crossings of $f_i=1$.
\end{proof}

The point of the lemma is that $\eta(f_i)$ need not be a rational
multiple of a single logarithm. At the final specialization $t=1/2$,
the needed values are independent linear forms in $A=\log2$ and
$B=\log3$. Thus the one-logarithm golden descent does not have to be
applied outside its stated hypotheses.

\subsection{The complete finite certificate}

For the following forty-two rows define
\begin{equation}
 f_i(t)=\varepsilon_i2^{d_i}t^{a_i}(1-t)^{b_i}(1+t)^{c_i}.
\label{rattetra:eq:arguments}
\end{equation}
All $a_i\ge0$, and all factors except the displayed sign are positive
for $0<t<1$. The coefficients below are primitive integers.

\begin{center}\small
\setlength{\tabcolsep}{4pt}
\begin{tabular}{r r c r r r r@{\qquad}r r c r r r r}
\toprule
$i$&$v_i$&$\varepsilon_i$&$d_i$&$a_i$&$b_i$&$c_i$&
$i$&$v_i$&$\varepsilon_i$&$d_i$&$a_i$&$b_i$&$c_i$\\
\midrule
1 & -1 & + & -1 & 0 & 0 & 1 & 22 & 3 & + & -1 & 1 & 0 & -1 \\
2 & -11 & + & 0 & 0 & 0 & 1 & 23 & -4 & + & 0 & 1 & 0 & -1 \\
3 & -6 & - & 0 & 0 & 0 & 1 & 24 & 6 & - & 0 & 1 & 0 & -1 \\
4 & -3 & + & 1 & 0 & 0 & 1 & 25 & 1 & + & 1 & 1 & 0 & -1 \\
5 & -4 & + & 0 & 0 & 1 & -1 & 26 & -3 & + & -1 & 1 & 0 & 0 \\
6 & 4 & - & 0 & 0 & 1 & -1 & 27 & -3 & - & -1 & 1 & 0 & 0 \\
7 & -1 & + & -1 & 0 & 1 & 0 & 28 & 6 & + & 0 & 1 & 0 & 0 \\
8 & 22 & + & 0 & 0 & 1 & 0 & 29 & -6 & - & 0 & 1 & 0 & 0 \\
9 & 21 & - & 0 & 0 & 1 & 0 & 30 & 3 & + & 1 & 1 & 0 & 0 \\
10 & -3 & + & 1 & 0 & 1 & 0 & 31 & 3 & - & 1 & 1 & 0 & 0 \\
11 & 4 & + & 0 & 0 & 1 & 1 & 32 & 1 & + & -1 & 1 & 0 & 1 \\
12 & -1 & + & 0 & 0 & 2 & 1 & 33 & 1 & - & -1 & 1 & 1 & 0 \\
13 & -1 & - & 0 & 0 & 3 & 0 & 34 & 3 & + & 0 & 1 & 1 & 0 \\
14 & -1 & + & -1 & 1 & -2 & 0 & 35 & 3 & - & 0 & 1 & 1 & 0 \\
15 & -6 & - & 0 & 1 & -2 & 0 & 36 & -4 & - & 0 & 2 & -1 & -1 \\
16 & -6 & - & 0 & 1 & -1 & -1 & 37 & -3 & - & 0 & 2 & -1 & 0 \\
17 & 3 & - & -1 & 1 & -1 & 0 & 38 & -1 & + & 1 & 2 & -1 & 0 \\
18 & -21 & + & 0 & 1 & -1 & 0 & 39 & -3 & + & 0 & 2 & 0 & -1 \\
19 & -7 & - & 0 & 1 & -1 & 0 & 40 & -1 & + & 1 & 2 & 0 & -1 \\
20 & 1 & - & 1 & 1 & -1 & 0 & 41 & 1 & + & 0 & 3 & -3 & 0 \\
21 & -1 & - & -1 & 1 & 0 & -2 & 42 & 1 & - & 0 & 3 & -1 & -2 \\
\bottomrule
\end{tabular}
\end{center}

\begin{lemma}[Exact tensor cancellation]
\label{rattetra:lem:certificate}
The displayed rational functions satisfy
\begin{equation}
 \sum_{i=1}^{42}v_i\beta_4[f_i]=0.
\label{rattetra:eq:tensor}
\end{equation}
\end{lemma}

\begin{proof}
The complete factor list needed for all $f_i$ and $1-f_i$, modulo
sign, is
\[
 2,\quad t,\quad t-1,\quad t+1,\quad
 t-2,\quad t+2,\quad2t-1,\quad2t+1,\quad
 t^2-t-1,\quad t^2-t+1.
\]
Use their exponent vectors as coordinates in $V$. For example,
the constant $2$ occurs in
\[
 1-\frac{1+t}{2}=\frac{1-t}{2},\qquad
 1-\frac{t}{2(1-t)^2}
   =\frac{(2t-1)(t-2)}{2(1-t)^2}.
\]
The first two slots of $\beta_4$ are symmetric. For each of the ten
unordered pairs of coordinates of $2,t,1-t,1+t$, multiply the
corresponding exponents and the two-by-two alternating minors of
the vectors of $f_i,1-f_i$, and sum with coefficient $v_i$.
Every resulting coordinate is zero. There are $230$ coordinates
touched by the finite sum.

The accompanying \texttt{tetralogarithm\_certificate.json} gives
exactly the forty-two displayed rows. The self-contained verifier
independently factors every rational function, checks the polynomial
factorizations by multiplication, retains the prime-$2$ coordinate,
and performs the indicated sums over rational integers. It does not
search for a relation or use numerical polylogarithms in this check.
This is a finite exact verification of~\eqref{rattetra:eq:tensor}.
\end{proof}

\subsection{The functional identity proved by the certificate}

For $0<t<1$, put
\[
 A=\log2,\qquad X=\log t,\qquad
 Y=\log(1-t),\qquad Z=\log(1+t),
\]
and define
\begin{align}
P_4={}&\frac{A^4}{4}-\frac{A^3(Y+Z)}6
       +\frac{A^2(Y^2+Z^2)}4+\frac{5A(Y^3+Z^3)}6
       -XY(3Y+Z)(Y+Z)\notag\\
 &+\frac{17Y^4+8Y^3Z+9Y^2Z^2+2YZ^3+2Z^4}{6},
 \label{rattetra:eq:P4}\\
P_2={}&-6A^2-4A(Y+Z)+\frac52Y^2+8YZ-8Z^2.
 \label{rattetra:eq:P2}
\end{align}

\begin{theorem}[A rational-function identity]
\label{rattetra:thm:functional}
With the forty-two rows of the table,
\begin{equation}
 \sum_{i=1}^{42}v_i\operatorname{Re}\operatorname{Li}_4(f_i(t))
 -4\operatorname{Li}_4(\tfrac12)
 =P_4+\zeta(2)P_2-\frac{71}{4}\zeta(4),
 \qquad 0<t<1.
\label{rattetra:eq:functional}
\end{equation}
The principal real parts and continuous values are used when an
argument is greater than one or equals one.
\end{theorem}

\begin{proof}
Fix $t\in(0,1)$ and write $L_i=\log|f_i(t)|$.
Apply Lemma~\ref{rattetra:lem:descent} with the functional sending a
prime or irreducible factor $2,s,1-s,1+s$ respectively to
$A,X,Y,Z$, and extend it linearly to the factor space, assigning
arbitrary values to its other basis elements. Unique factorization
shows that such a functional exists, and its value on every
$f_i(s)$ is $L_i$. This definition remains valid when one of the
additional factors of $1-f_i(t)$ vanishes; no logarithm of zero
is used to define the functional. The table functions have limits
\[
 e_i=\lim_{s\to0^+}f_i(s)=
 \begin{cases}0,&a_i>0,\\ \varepsilon_i2^{d_i},&a_i=0.
 \end{cases}
\]
Thus for $n=2,3,4$,
\[
 \sum_i v_iL_i^{4-n}R_n(f_i(t))
  =\sum_{a_i=0}v_iL_i^{4-n}R_n(e_i).
\]
The elementary triangular identity
\[
 \ell_4(x)=R_4(x)+L R_3(x)+\tfrac12L^2R_2(x)
                         +\tfrac1{24}L^3\ell_1(x)
\]
then supplies an exact ordinary-polylogarithm identity. More explicitly,
put $M_i=\log|e_i|=d_iA$ for $a_i=0$, and set
$Q_i=-M_i^3/8+L_iM_i^2/3-L_i^2M_i/4$. The result is
\begin{align}
\sum_i v_i\ell_4(f_i(t))={}&\sum_{a_i=0}v_i\bigl\{
 \ell_4(e_i)+(L_i-M_i)\ell_3(e_i)
 +\tfrac12(L_i-M_i)^2\ell_2(e_i)+Q_i\ell_1(e_i)\bigr\}\notag\\
 &+\frac1{24}\sum_i v_iL_i^3\ell_1(f_i(t)).
\label{rattetra:eq:descent-master}
\end{align}
When $e_i=1$, $Q_i=0$ and the corresponding last term in braces
is zero; no divergent value is assigned to $\ell_1(1)$.
Likewise $L_i^3\ell_1(f_i(t))$ is assigned its limiting value zero
at $f_i(t)=1$.

Only $e_i\in\{1,-1,1/2,2\}$ occur. The required lower values are
\[
\begin{array}{c|ccc}
e&\ell_1(e)&\ell_2(e)&\ell_3(e)\\ \hline
-1&-A&-\zeta(2)/2&-3\zeta(3)/4\\
1/2&A&\zeta(2)/2-A^2/2&7\zeta(3)/8-\zeta(2)A/2+A^3/6\\
2&0&3\zeta(2)/2&7\zeta(3)/8+3\zeta(2)A/2
\end{array}
\]
and $\ell_2(1)=\zeta(2)$, $\ell_3(1)=\zeta(3)$.
For completeness, the lower reflection and Landen identities are
\begin{align*}
\operatorname{Li}_2(x)+\operatorname{Li}_2(1-x)
 &=\zeta(2)-\log x\log(1-x),\\
\operatorname{Li}_2\!\left(-\frac{x}{1-x}\right)
 &=-\operatorname{Li}_2(x)-\frac12\log^2(1-x),\\
\operatorname{Li}_3(x)+\operatorname{Li}_3(1-x)
 +\operatorname{Li}_3\!\left(-\frac{x}{1-x}\right)
 &=\zeta(3)+\zeta(2)\log(1-x)\\
 &\quad-\frac12\log x\log^2(1-x)+\frac16\log^3(1-x),
\end{align*}
for $0<x<1$. Differentiating the first two identities gives zero
differences, with constants fixed at $x\to0^+$. Using them in the
derivative of the third gives zero as well, and the same endpoint
fixes its constant. Substitution $x=1/2$ proves the displayed half
values. The values at $2$ follow by real inversion. Thus these
endpoint reductions require no integer-relation recognition.

The endpoint weight-four sum is exactly
\[
 \sum_{a_i=0}v_i\ell_4(e_i)
 =4\operatorname{Li}_4(\tfrac12)+\frac{A^4}{4}
       -6\zeta(2)A^2-\frac{71}{4}\zeta(4).
\]
Substitute $L_i=d_iA+a_iX+b_iY+c_iZ$ and the factorizations of
$1-f_i$ in~\eqref{rattetra:eq:descent-master}. All logarithms of the
six additional polynomial factors cancel, as do the $\zeta(3)$ terms.
The remaining polynomial is precisely
$P_4+\zeta(2)P_2-71\zeta(4)/4$, which proves the theorem.
The verifier expands both the master formula and the displayed
$P_4,P_2$ independently and checks their exact polynomial difference
is zero.
\end{proof}

\subsection{The rational specialization, including all inversion terms}

For $r>1$, the real inversion formulas at weight four are
\begin{align*}
 \ell_4(r)+\operatorname{Li}_4(r^{-1})
   &=-\frac{\log^4r}{24}+\zeta(2)\log^2r+2\zeta(4),\\
 \operatorname{Li}_4(-r)+\operatorname{Li}_4(-r^{-1})
   &=-\frac{\log^4r}{24}-\frac12\zeta(2)\log^2r
       -\frac74\zeta(4).
\end{align*}
They follow by repeated differentiation from the real logarithm
and dilogarithm inversion laws, fixing the constants at $r=1$.
Negative arguments inside the unit interval are removed by
\[
 \operatorname{Li}_4(-r)=\tfrac18\operatorname{Li}_4(r^2)
                                  -\operatorname{Li}_4(r),
 \qquad 0\le r\le1.
\]
At $t=1/2$ one has $X=Y=-A$, $Z=B-A$ and
$f_i(1/2)=\varepsilon_i2^{d_i-a_i-b_i-c_i}3^{c_i}$.
Applying the three displayed rules to the left side of
\eqref{rattetra:eq:descent-master}, after subtracting its endpoint
weight-four sum and multiplying by $180$, gives exactly the six
tetralogarithms in~\eqref{rattetra:eq:main}. The remaining polynomial
from those inversions is
\begin{align*}
I={}&120A^4-510A^3B+765A^2B^2-510AB^3+150B^4\notag\\
 &+\zeta(2)(90A^2+2880AB-1980B^2)-1710\zeta(4).
\end{align*}
The non-weight-four part of the specialized descent master formula is
\begin{align*}
D={}&-450A^4+210A^3B+225A^2B^2-150AB^3+60B^4\notag\\
 &+\zeta(2)(1890A^2+720AB-1440B^2).
\end{align*}
Consequently the six-term sum equals $D-I$, namely
\begin{align*}
&-570A^4+720A^3B-540A^2B^2+360AB^3-90B^4\notag\\
&\qquad+\zeta(2)(1800A^2-2160AB+540B^2)+1710\zeta(4).
\end{align*}
Using $\zeta(2)=\pi^2/6$, $\zeta(4)=\pi^4/90$ gives exactly
\eqref{rattetra:eq:main}, proving Theorem~\ref{rattetra:thm:main}.

% Optional insertion after identities_notes.tex; the certificate is shared.
% Requires the same theorem environments, amsmath and booktabs.
\subsection{A further proved rational specialization}
\label{rattetra:sec:onethird}

The functional identity also produces explicit identities beyond the
six-term conjecture. The specialization $t=1/3$ is convenient because
all logarithms reduce to $A=\log2$ and $B=\log3$. The following
corollary is a new entry in this report, with no claim of priority
among classical tetralogarithm identities.

\begin{corollary}[Twenty-five positive rational arguments]
\label{rattetra:cor:onethird}
For the argument--coefficient pairs $(r,c_r)$ in the table,
\begin{align}
\sum_r c_r\operatorname{Li}_4(r)
={}&\frac{109}{3}A^4-116A^3B+54A^2B^2+36AB^3-21B^4\notag\\
 &+\zeta(2)(316A^2-408AB+132B^2)+53\zeta(4).
\label{rattetra:eq:onethird}
\end{align}
All arguments lie in $(0,1)$, and the integer coefficient list is
primitive: its greatest common divisor is one.

\begin{center}
\renewcommand{\arraystretch}{1.25}
\begin{tabular}{c r@{\qquad\qquad}c r}
\toprule
$r$&$c_r$&$r$&$c_r$\\
\midrule
$1/1024$ & $1$ & $9/64$ & $-6$ \\
$9/1024$ & $-1$ & $1/6$ & $16$ \\
$1/81$ & $1$ & $2/9$ & $8$ \\
$1/64$ & $-4$ & $1/4$ & $-107$ \\
$1/36$ & $-6$ & $8/27$ & $8$ \\
$1/32$ & $-8$ & $1/3$ & $80$ \\
$4/81$ & $3$ & $3/8$ & $64$ \\
$1/16$ & $9$ & $4/9$ & $24$ \\
$1/12$ & $-24$ & $1/2$ & $-200$ \\
$64/729$ & $-1$ & $16/27$ & $-8$ \\
$3/32$ & $8$ & $3/4$ & $112$ \\
$1/9$ & $-14$ & $8/9$ & $32$ \\
$1/8$ & $64$ & & \\
\bottomrule
\end{tabular}
\end{center}
\end{corollary}

\begin{proof}
Set $t=1/3$ in Theorem~\ref{rattetra:thm:functional}. The logarithms
are $X=-B$, $Y=A-B$, and $Z=2A-B$. Apply the real inversion
formulas to arguments whose absolute value exceeds one, then the
duplication formula
$\operatorname{Li}_4(-r)=\operatorname{Li}_4(r^2)/8-
\operatorname{Li}_4(r)$ to negative arguments in $[-1,0]$.
Multiplication by eight clears the resulting coefficient denominators.
Collecting equal positive arguments gives exactly the table, and
collecting all inversion terms gives the displayed polynomial.
The same certificate verifier performs these reductions over
$\mathbb Q$ and checks that the ordinary-polynomial residual is zero.
As an independent audit only, a direct evaluation of the reduced
identity at ninety working decimal digits has absolute residual
below $8\times10^{-90}$.
\end{proof}

The argument list and polynomial are supplied separately in
\src{tetralogarithm_one_third.json}, and are also included in the
shared certificate data. The replay checks the original conjecture,
the functional identity, and this corollary in a single run.

\subsection{Verification, provenance, and further questions}

The exact verifier finishes with a zero tensor residual, a zero
ordinary-polynomial residual, and the six specified rational arguments
with coefficients $(720,-2160,2160,270,540,135)$. Separately, it evaluates
the functional identity at $t=1/5,37/100,1/2,4/5$ using $90$ working
decimal digits. The largest numerical discrepancy is approximately
$7.9\times10^{-90}$. These calculations check the branch conventions;
the proof is the finite tensor cancellation and the analytic descent.

The original questions are
\url{https://math.stackexchange.com/questions/1373123/} and
\url{https://mathoverflow.net/questions/261408/}.
The former still presents a proposed Kummer route, while the latter's
answer identifies the separate golden weight-five ladders. Those two
facts should not be conflated: the certificate here settles the rational
weight-four formula directly. The method is compatible with the
higher-Bloch-condition approach to functional equations, but no unproved
period conjecture or transcendence assertion is used.

Three useful next problems are to decompose the forty-two rows into a
small set of named classical Kummer specializations; to specialize the
proved functional identity at other rational or quadratic arguments and
systematically reduce the resulting baskets; and to search for similarly
complete certificates for the unresolved golden weights six and seven.
The Gaussian $S_6$ conjecture remains a distinct problem and is not
promoted by this result. No minimality or arithmetic independence of
the six rational tetralogarithms is asserted.
